\documentclass[conference]{IEEEtran}
\IEEEoverridecommandlockouts
% The preceding line is only needed to identify funding in the first footnote. If that is unneeded, please comment it out.
\usepackage{cite}
\usepackage{amsmath,amssymb,amsfonts}
\usepackage{algorithmic}
\usepackage{graphicx}
\usepackage{textcomp}
\usepackage{xcolor}
\def\BibTeX{{\rm B\kern-.05em{\sc i\kern-.025em b}\kern-.08em
    T\kern-.1667em\lower.7ex\hbox{E}\kern-.125emX}}
\begin{document}

\title{Unifying Isolated Spin-Down and Binary Recycling in Massive Neutron Star Formation via Relaxed Equation of State Axioms}\author{Anonymous Research Plan Author}\maketitle

\begin{abstract}
This proposal advances a candidate mechanism for massive neutron star formation by relaxing static equilibrium axioms of the dense matter equation of state to incorporate rotational and thermal support from binary mass accretion. The central contribution is the formal unification of isolated spin-down and binary recycling pathways, predicting a distinct mass-spin relation that differentiates massive systems exceeding the static Tolman-Oppenheimer-Volkoff limit from standard isolated magnetic-dipole spin-down pulsars. Existing observational constraints on neutron star radii and equation of state parameters, particularly those derived from GW170817, establish the necessary causal bounds for this theoretical model. The argument integrates evidence from binary evolution simulations and pulsar timing data to define the validity domain where accretion-driven angular momentum transfer modifies effective equilibrium conditions without violating causality. This route specifically addresses the gap in mechanistic explanation for massive millisecond pulsars by proposing that companion ablation primarily affects surface and magnetospheric properties, allowing the core mass-spin relation to remain governed by the relaxed axioms.

The proposed theoretical framework requires verification of specific proof obligations to establish the mass-spin relation as a valid predictor for systems exceeding 2.0 solar masses. The research design maps prespecified outcome branches to these formal dependencies, distinguishing between conditions that support the relaxed axioms and those that favor standard rotational support models. A null or contradictory outcome is defined as a failure of the relaxed axioms to maintain causality or a prediction that contradicts the observed population of massive millisecond pulsars. Conversely, an uninformative branch arises when validity criteria prevent interpretation, necessitating human review of measurement and sampling prerequisites. The argument transitions to the derivation stage by establishing that the distinct mass-spin relation serves as the primary discriminator between evolutionary channels, contingent upon the successful resolution of formal reasoning gaps regarding the exact mathematical relaxation conditions.
\end{abstract}
\begin{IEEEkeywords}
neutron stars, equation of state, binary recycling, mass-spin relation, rotational support, TOV limit
\end{IEEEkeywords}\section{Introduction}
\subsection{The Observational-Theory Divide}
Neutron star astrophysics currently operates across two largely disconnected formal regimes: static Tolman-Oppenheimer-Volkoff (TOV) equilibrium constrained by causal equations of state, and dynamical binary evolution governed by mass transfer and angular momentum accretion. The static TOV limit defines the maximum mass for non-rotating, cold neutron stars, while binary recycling pathways explain the existence of millisecond pulsars through accretion-driven spin-up. A persistent mechanism\_explanation\_gap exists because these two regimes are rarely unified under a single formal framework. Current models treat the equation of state as a static input to binary evolution codes, rather than a dynamic variable subject to thermodynamic and mechanical relaxation. This separation obscures whether the observed high masses of millisecond pulsars require fundamental modifications to dense matter physics or are fully accounted for by standard rotational support.~\cite{cite_55957f4ce67bd972,cite_7e1f0038b29c01f3}

\subsection{Relaxed Equilibrium Axioms as a Candidate Mechanism}
This proposal introduces a candidate mechanism that relaxes the static equilibrium axioms of the dense matter equation of state to incorporate rotational and thermal support from binary mass accretion. Under this framework, the effective equilibrium conditions of the neutron star core are modified by continuous thermodynamic and mechanical perturbations. The central hypothesis asserts that this relaxation permits masses beyond the static TOV limit without violating the causality bound, defined as the speed of sound remaining less than or equal to the speed of light. By formally integrating binary accretion parameters into the equilibrium axioms, the model predicts a distinct mass-spin relation. This relation serves as a theoretical bridge, unifying isolated spin-down and binary evolution pathways and offering a discriminative observable for massive neutron stars exceeding 2.0 solar masses.

\subsection{Scope and Falsification Boundaries}
The contribution is scoped as a candidate\_mechanism subject to rigorous formal and observational boundaries. The proposal fails if the relaxed static equilibrium axioms inevitably violate the causality bound, if the predicted mass-spin relation contradicts the observed population of massive millisecond pulsars, or if standard rotational support models fully explain the observations without requiring axiom relaxation. A key alternative explanation remains that high masses are solely due to binary spin-up acting on a standard stiff equation of state. Distinguishing this proposal from the alternative requires verifying that the derived mass-spin relation yields predictions distinct from isolated magnetic-dipole spin-down tracks. The work is theoretical and static; it does not execute new simulations or observations but establishes the formal conditions and proof obligations necessary to test the mechanism.~\cite{cite_6c11517977d0b39d}

\subsection{Transition to Formal Verification}
The argument advances from the identification of the mechanism gap to the specification of the formal verification plan. The next stage details the definition ledger, including the mathematical relaxation conditions and the specific functional form of the mass-spin relation. It establishes the proof obligations required to demonstrate consistency with the causality bound and defines the decision matrix for prespecified outcome branches. This transition moves the proposal from a conceptual unification strategy to a concrete, auditable formal reasoning plan.

\section{Background, Survey, and Research Gap}
\subsection{Dense-Matter Constraints and the Static TOV Limit}
The maximum gravitational mass of a neutron star is set by the dense-matter equation of state (EoS) under the assumption of static, non-rotating equilibrium. Unified EoS treatments anchor the crust to ground-state composition and the liquid core to a minimal npe$\mu$ composition, providing the structural baseline for Tolman-Oppenheimer-Volkoff (TOV) mass limits (W2103740830). Observational constraints on the radius and maximum mass are not independent of this structural assumption: the inferred radius bounds from the binary neutron star merger GW170817 depend on whether the analysis additionally requires the EoS to support neutron stars with masses larger than 1.97 solar masses, a threshold indicated by electromagnetic observations (W2807309204). Joint analyses that incorporate heavy-ion collision data into astrophysical constraints further improve the pressure bound at supranuclear densities (W3182570346). These results establish the static limit as a well-defined, observationally testable boundary, against which any proposed relaxation of the equilibrium axioms must be measured.~\cite{cite_42e341ff6471ce89,cite_55957f4ce67bd972,cite_6c11517977d0b39d}

\subsection{Binary Recycling and the Mass-Spin Population}
Massive neutron stars are produced through binary evolution, where mass transfer and angular momentum exchange spin up the remnant to millisecond periods. Binary interaction via winds, tides, mass transfer, and mergers directly modulates the rotation-rate distribution of massive stars (W2001595193). The accretion torque that drives recycling is counterbalanced by wind-induced torques and magnetic braking, whose efficiencies are bounded by observational constraints on the mass outflow rate (W2012578543). The resulting population of millisecond pulsars is characterized by precisely measured astrometric and timing parameters derived from long-baseline radio campaigns (W2181312162). Within this population, the black widow pulsar PSR J0952-0607 exemplifies the extreme endpoint of ablation-driven evolution, with a measured mass of 2.35 $\pm$ 0.17 M\_sun that places a stringent lower bound on the maximum neutron star mass independent of Shapiro-delay systems. The existence of such massive, rapidly rotating remnants is an established observational fact; the unresolved question is the physical mechanism that permits these masses to coexist with millisecond periods without violating the static EoS limit.~\cite{cite_1e5d552669fc0c56,cite_2908f4a456e2afd6,cite_be7a68ae6d0cde0b,cite_8609757b321d8ce0,cite_fd6df2b02dbddc60}

\subsection{The Causal Boundary and the Mechanism Gap}
Any relaxation of the static equilibrium axioms to accommodate masses beyond the static TOV limit must preserve the causal speed bound, defined as the condition that the speed of sound in dense matter remains below the speed of light. The proposal under consideration treats binary mass accretion and angular momentum transfer as a continuous thermodynamic and mechanical perturbation that modifies the effective equilibrium conditions, permitting mass accumulation beyond the static limit while preserving causality. This candidate mechanism is not yet covered in the relevant evidence cluster on massive neutron star masses and dense matter equations of state. The gap is structural: the observed relation between mass and spin in the massive millisecond population lacks an identified mechanism that can distinguish it from competing explanations. The survey establishes the boundary conditions---causality, EoS composition, and binary angular momentum transfer---within which this gap must be resolved, but it does not supply the formal relaxation conditions themselves.~\cite{cite_42e341ff6471ce89,cite_55957f4ce67bd972}

\subsection{Alternative Explanation: Standard Stiff EoS with Rotational Support}
The high masses of millisecond pulsars are explained by standard stiff EoS combined with rotational support exceeding the static TOV limit due to binary spin-up, without requiring any relaxation of the fundamental nuclear EoS axioms. This alternative is fully consistent with the established observational record: the measured masses and periods are accounted for by angular momentum accretion and centrifugal support within the static framework. The competing explanations are therefore not distinguished by the existence of massive millisecond pulsars alone. The discriminative power of the relaxed-axiom proposal rests entirely on whether it predicts a mass-spin relation that is distinct from the relation produced by standard rotational support. The survey evidence does not yet supply a formal comparison between these two relations; the gap is precisely the absence of a derived, falsifiable relation under the relaxed axioms.~\cite{cite_1e5d552669fc0c56,cite_55957f4ce67bd972}

\subsection{Missing Premises and Boundary Variables}
\begin{itemize}
\item ['The specific mathematical form of the relaxation applied to the static equilibrium axioms is referenced but not defined. Without a formal statement of the relaxed axioms, the candidate mechanism remains a proposal rather than a derivable relation.', 'The functional form of the mechanism mass-spin relation, denoted as M = R(P), is referenced but not specified. The derived relation that would distinguish the relaxed-axiom population from the isolated spin-down population is the central unverified object.', 'The equations governing binary mass and angular momentum transfer require formal specification to be treated as a continuous perturbation to the equilibrium conditions rather than a discrete accretion event.', 'The formal criteria for testing causality violation under the relaxed axioms are missing. The causality bound is declared as a control variable, but the test that would confirm its preservation under relaxation is not yet defined.', 'The measurement method for masses exceeding 2.0 solar masses is not specified in the survey evidence. Whether Shapiro delay or another timing observable anchors the mass measurement determines the observational validity of the comparison population.', 'The measurement method for the spin period under the relaxed-axiom prediction is not specified. The period is the independent variable in the proposed relation, and its operational definition must be fixed before the relation can be tested.']
\end{itemize}~\cite{cite_1e5d552669fc0c56,cite_55957f4ce67bd972}

\subsection{Transition to the Formal Argument}
The survey evidence establishes the static EoS limit, the binary recycling channel, and the causal bound as the fixed boundary conditions of the problem. The gap is the absence of a formal derivation of the relaxed equilibrium axioms and the resulting mass-spin relation. The next section must define the candidate theorem's domain, state the admissible premises, and supply the scoped entry lemma that formalizes the relaxation. The comparison between the relaxed-axiom relation and the standard rotational-support relation is the decisive test; it is not resolved by the existing literature and must be constructed as a formal proof obligation.

\section{Research Questions and Planned Contributions}
\subsection{Operational Research Questions}
\begin{itemize}
\item **RQ1:** What exact mathematical relaxation conditions permit the static equilibrium axioms of the dense matter equation of state to support masses beyond the static Tolman-Oppenheimer-Volkoff (TOV) limit via rotational and thermal support without violating the causality bound?
\item **RQ2:** Does the resulting candidate mass-spin relation predict a distinct evolutionary track for massive neutron stars (>2.0 M\_sun) with millisecond periods that is formally distinguishable from isolated magnetic-dipole spin-down pulsars?
\item **RQ3:** Which numerical methods or analytical techniques will verify the proposed mass-spin relation and resolve the discarded counterexample analysis?
\end{itemize}~\cite{cite_42e341ff6471ce89,cite_a50fdc75b0400b49}

\subsection{Planned Contributions and Argumentative Roles}
The research plan operationalizes a candidate mechanism for massive neutron star formation. The first contribution is a scoped entry lemma formalizing how binary accretion modifies effective equilibrium conditions, establishing the domain where mass accumulation exceeds the static TOV limit while preserving the speed-of-sound constraint. The second contribution is a dependency-closure matrix linking the relaxed axioms to a proposed mass-spin relation, distinguishing binary recycling from isolated spin-down tracks. The third contribution is a prespecified decision protocol mapping outcome branches to proof obligations, ensuring that formal claims remain separate from empirical constraints. These contributions address the mechanism explanation gap in dense matter equations of state by replacing standard static axioms with a continuous thermodynamic and mechanical perturbation model.~\cite{cite_55957f4ce67bd972,cite_1e5d552669fc0c56,cite_c369df652181e8f7}

\subsection{Unresolved Items Delimiting the Questions}
The scope of these questions is delimited by three unresolved design dependencies. The exact mathematical relaxation conditions for the static equilibrium axioms remain undefined, requiring formal specification before the entry lemma can be verified. The specific numerical methods or analytical techniques for verifying the mass-spin relation are unspecified, preventing the execution of the proposed decision protocol. Additionally, the discarded counterexample analysis batch requires qualified human review to validate the boundary conditions and confirm that the relaxed axioms do not inevitably violate causality or tidal deformability constraints. These dependencies define the boundary of the current proposal and establish the prerequisites for the next stage of formal derivation.~\cite{cite_42e341ff6471ce89,cite_a50fdc75b0400b49}

\section{Idea Source Checkpoints and Direction Selection Audit}
The selected direction, mechanism\_replacement, originates from a survey-identified gap: the required slot 'candidate\_mechanism' remains uncovered in the evidence cluster concerning Massive Neutron Star Masses and Dense Matter Equations of State. This gap highlights the absence of a causal pathway connecting collapse inputs to a magnetized rotating neutron star that distinguishes successful pulsar formation from failed collapses. The proposal addresses this by relaxing static equilibrium axioms of the dense matter equation of state (EoS) to incorporate rotational and thermal support from binary mass accretion. This theoretical intervention aims to unify isolated spin-down and binary recycling, predicting a distinct mass-spin relation for neutron stars exceeding the static Tolman-Oppenheimer-Volkoff (TOV) limit. The scientific attraction lies in formalizing the transition from phenomenological descriptions to predictive models where centrifugal support permits mass accumulation beyond static limits without violating causality bounds.

\begin{itemize}
\item The idea evolution audit reveals three available checkpoints, though temporal ordering is unknown. Checkpoint 1 (idea\_candidate.json) establishes the core hypothesis of relaxing static EoS axioms for massive neutron star formation. Checkpoint 2 (idea\_portfolio.json) retains the central hypothesis but shifts the title focus toward dynamical inclination and wind braking in isolated pulsar spin-down, indicating a potential divergence in emphasis between binary recycling mechanisms and isolated spin-down physics. Checkpoint 3 (idea\_result.json) consolidates the direction into a mechanism\_replacement mode, explicitly listing assumptions about continuous thermodynamic perturbation via binary accretion and the primary effect of companion ablation on surface/magnetosphere properties. These snapshots confirm the direction's stability around the candidate\_mechanism scope but expose a lack of formal definition for the specific mathematical relaxation conditions.
\end{itemize}

Critical defects in the upstream reasoning constrain the current proposal's strength. The counterexample\_analyzer batch was discarded due to invalid response, leaving no generated formalization or counterexample for the relaxed axioms. Consequently, the specific functional form of the 'mechanism\_mass\_spin\_relation' remains undefined, and the mathematical relaxation conditions for the static equilibrium axioms lack formal specification. The proposal assumes that binary accretion can be modeled as a continuous perturbation modifying effective equilibrium conditions, but this link is unverified. Furthermore, the alternative explanation that standard stiff EoS combined with rotational support fully explains massive millisecond pulsars without relaxing fundamental nuclear EoS axioms remains a live competing mechanism. The current design does not yet provide a discriminative criterion to rule out this standard model, relying instead on the unverified prediction that the relaxed axioms yield a distinct mass-spin relation.

The retained direction is qualified as a candidate\_mechanism requiring human review before formal proof obligations are engaged. The proposal explicitly excludes the modeling of common envelope and contact binary phases, a boundary condition supported by MESA binary evolution module limitations. The core contribution is the formal relaxation of static equilibrium axioms to permit masses beyond the static TOV limit via rotational support, subject to the causality bound (c\_s <= c). This direction is selected over the isolated spin-down focus of Checkpoint 2 because it directly addresses the mechanism gap in massive neutron star formation. However, the transition to the next stage depends on resolving the undefined mathematical relaxation conditions and establishing a verification plan for the causality bound. The proposal remains at the design assumption stage, with no observed results or completed proofs.~\cite{cite_c369df652181e8f7}

\section{Problem Definition, Assumptions, and Hypotheses}
This section establishes the formal problem for a candidate mechanism linking massive neutron star formation to binary recycling. The central proposal is that relaxing the static equilibrium axioms of the dense matter equation of state (EoS) to incorporate rotational and thermal support permits masses beyond the static Tolman-Oppenheimer-Volkoff (TOV) limit without violating the causality bound. This relaxation is posited to unify isolated spin-down and binary evolution pathways, predicting a distinct mass-spin relation for massive neutron stars. The argument consumes the candidate theorem entry, specifically the assumptions and propositions regarding the modified equilibrium conditions.

\paragraph{Definition.} The core formal object is the relaxed equilibrium relation M = R(P), where M represents the gravitational mass of the neutron star and P represents the spin period. The function R(P) is derived from the relaxed static equilibrium axioms. The domain of this relation is constrained by the causality bound, defined as the condition that the speed of sound in dense matter remains below the speed of light. The admissible premises for this definition include the continuous thermodynamic and mechanical perturbation induced by binary accretion support.

\begin{equation}
c_s <= c
\label{eq:formal-problem-and-hypotheses-b3}
\end{equation}

\begin{itemize}
\item The assumption ledger for the candidate theorem entry comprises three primary premises. First, the static equilibrium axioms of the dense matter EoS can be formally relaxed to permit mass accumulation beyond the static TOV limit via rotational and thermal support without violating the causality bound. Second, binary mass accretion and angular momentum transfer can be modeled as a continuous thermodynamic and mechanical perturbation that modifies the effective equilibrium conditions of the neutron star. Third, companion ablation primarily affects the surface and magnetosphere, allowing the core mass-spin relation to remain governed by the relaxed effective equilibrium axioms.
\end{itemize}

\paragraph{Lemma Registry L1, L2 (Candidate).} Lemma L1 (Candidate): Under the assumptions of causal EoS bounds, relaxing the static equilibrium axioms to incorporate rotational support from binary mass accretion permits masses beyond the static TOV limit. Lemma L2 (Candidate): The relaxed axioms predict a distinct mass-spin relation distinguishing massive neutron stars (>2.0 M\_sun) with millisecond periods from isolated magnetic-dipole spin-down pulsars. These lemmas are proposed and unverified, requiring formal proof obligations to establish their validity within the defined domain.

\paragraph{Proof Obligation Registry PO1, PO2 (Unverified).} The primary proposition is that the relaxation of equilibrium axioms allows for masses beyond the static TOV limit, while the derived relation distinguishes millisecond massive stars. The failure condition for this proposition is that the relaxed static equilibrium axioms violate the causality bound, or the predicted mass-spin relation contradicts the observed population of massive millisecond pulsars, or standard rotational support models fully explain the observations without needing the relaxed axioms. Proof obligations PO1 and PO2 are required to verify the consistency of the relaxed axioms with the causality bound and the distinctiveness of the mass-spin relation.

\section{Study Design and Methods}
\subsection{Design Architecture and Operational Boundary}
This proposal adopts a static, theoretical-model design in which the unit of analysis is a candidate equation-of-state (EoS) framework rather than an observational sample. The central mechanism replaces the standard static equilibrium axioms of dense matter with relaxed axioms that formally admit rotational and thermal support from binary accretion. The design goal is not to measure a population but to specify a protocol that can derive and audit a candidate mass-spin relation, then test that relation against two declared alternatives: (i) standard stiff EoS plus ordinary rotational support exceeding the static Tolman-Oppenheimer-Volkoff (TOV) limit, and (ii) evolutionary channels fully explained without relaxing the static axioms. Formal claims and empirical claims remain separate; no observational result is produced here. The definition ledger (owner: definitions\_and\_propositions) supplies the symbols and assumptions consumed below, and the review ledger (owner: risk\_limitations\_and\_review) supplies the human-review gates; this section states only the dependency consequences, not their item-by-item contents.

\begin{table*}[htbp]
\centering
\footnotesize
\setlength{\tabcolsep}{3pt}
\renewcommand{\arraystretch}{1.12}
\caption{Comparator and Ablation Matrix}
\label{tab:study-design-and-methods-b2}
\begin{tabular}{|p{0.184\textwidth}|p{0.184\textwidth}|p{0.184\textwidth}|p{0.184\textwidth}|p{0.184\textwidth}|}
\hline
Design arm & Fixed inputs & Perturbed input & Decision criterion & Consequence if criterion fails \\\hline
Relaxed-axiom model & Causality bound c\_s <= c; core composition assumptions & Rotational/thermal support term & Derived relation R(P) yields M > M\_static\_max within causal bound & Revisit relaxation form; do not update theorem status \\\hline
Standard-support control & Same causality bound; same EoS family & No relaxation of static axioms & Ordinary spin-up reproduces observed massive MSPs without relaxation & Relaxation is not required; mechanism loses explanatory role \\\hline
Channel-separation control & Same R(P) candidate & Isolated magnetic-dipole spin-down track & Massive NSs (>2.0 M\_sun) with ms periods separate from isolated population & Discriminating claim unsupported in this boundary \\\hline
Numerical-implementation ablation & Same physics & Hydro/transfer discretization & Conclusions stable across resolution and scheme & Code-to-code variation dominates; comparison inconclusive \\\hline
\end{tabular}
\end{table*}

\subsection{Causality Control Constraint}
\begin{equation}
c_s^2 = dP/d\epsilon \leq c^2
\label{eq:study-design-and-methods-b3}
\end{equation}

\subsection{Boundary Checks and Comparison Adequacy}
The causality constraint is the load-bearing control variable: the relaxed axioms are admissible only if the resulting pressure-density curve keeps the squared sound speed at or below c\textasciicircum{}2 across the core. The design therefore treats the causality check as a hard boundary, not a tunable parameter. Comparison adequacy is judged against the declared alternative explanation that standard stiff EoS plus binary spin-up already accounts for massive millisecond pulsars; the relaxed-axiom model earns its place only if it produces a mass-spin relation that the standard model cannot reproduce within the same causal bound. Evidence from GW170817 shows that radius and EoS constraints are sensitive to whether the analysis additionally requires support for masses above 1.97 M\_sun, and that component-mass inference depends strongly on spin priors; these facts set the comparison's stakes and the prior structure the model must respect. Population-synthesis experience warns that omitting tidal evolution or mishandling initial orbital distributions introduces systematic bias, so the protocol fixes initial-condition distributions as controls rather than free parameters. Code-to-code sensitivity is an explicit confounder: because implementation choices for star formation, feedback, and mass-transfer schemes have been shown to drive large output variations even at fixed physics, the design mandates a resolution- and scheme-ablation arm before any comparison is declared informative. See Eq.~\eqref{eq:study-design-and-methods-b3}.~\cite{cite_3d2e6757608d5364,cite_42e341ff6471ce89,cite_81b31d0d8218f9b8}

\subsection{Audit Protocol and Human-Decision Gates}
\paragraph{Planned protocol.} The protocol proceeds in three prespecified stages, each mapped to a proof obligation or a no-information condition. Stage 1 fixes the control variables (causality bound, EoS equilibrium conditions, initial orbital distributions) and records the candidate relation R(P) as unverified formal work. Stage 2 runs the comparator and ablation arms; a branch is informative only if the causality check passes and the standard-support control fails to reproduce the same relation. Stage 3 evaluates channel separation against the isolated-spin-down track. The no-information branches are procedural, not scientific: if the relaxation form, the accretion-transfer equations, the functional form of R(P), or the causality-test criterion is unresolved, the dependency does not update theorem status and the design is neither validated nor falsified. Because the counterexample analysis stage was discarded upstream, no counterexample conclusion is drawn here; the design is treated as incomplete until qualified human review supplies the missing formal definitions and verification criteria. The next stage should consume this protocol as a checklist of open formal dependencies rather than as a source of results.~\cite{cite_58f0ac6612957738}

\section{Expected Outcome Branches and Conditional Conclusions}
\subsection{Decision Protocol Scope}
This section operationalizes the conditional decision logic for the candidate mechanism. The proposal maps prespecified outcome branches to the relevant Lemma and Proof Obligation, defines the allowed conclusion for each branch, and specifies the next action. No results are observed; the protocol governs how future formal or numerical verification would update the theorem status.

\begin{table*}[htbp]
\centering
\footnotesize
\setlength{\tabcolsep}{3pt}
\renewcommand{\arraystretch}{1.12}
\caption{Conditional Outcome Decision Matrix}
\label{tab:expected-outcomes-eob-matrix}
\begin{tabular}{|p{0.184\textwidth}|p{0.184\textwidth}|p{0.184\textwidth}|p{0.184\textwidth}|p{0.184\textwidth}|}
\hline
Outcome Branch & Trigger & Relevant Lemma / PO & Allowed Conclusion & Next Action \\\hline
supports\_mechanism & Prespecified analysis is consistent with the declared relation while planned controls do not favor a stated alternative explanation. & L1, L2, PO1, PO2 & Supports, but does not prove, the declared relation within the design boundary. & Replicate under independently confirmed conditions; test the most consequential declared boundary condition. \\\hline
partial\_or\_heterogeneous & Prespecified analysis indicates variation across declared conditions, units, or measurement contexts. & L1, L2, PO1, PO2 & The relation is conditional or heterogeneous; no universal conclusion is warranted. & Predefine and check plausible moderators; improve the coverage of conditions and measurement comparability. \\\hline
null\_or\_contradictory & Prespecified comparison does not support the declared relation or instead favors a declared alternative explanation. & L1, L2, PO1, PO2 & The proposed relation is not supported in this design boundary; absence of support is not proof of absence generally. & Audit construct validity and comparison adequacy; revise the mechanism or boundary claim before another design iteration. \\\hline
uninformative\_or\_invalid & Prespecified quality-control, missingness, protocol-deviation, or validity criteria prevent interpretation. & L1, L2, PO1, PO2 & No scientific conclusion is warranted because the planned design did not yield interpretable evidence. & Resolve the identified validity or data-quality failure before repeating the design; obtain human confirmation of measurement, sampling, and analysis prerequisites. \\\hline
\end{tabular}
\end{table*}

\subsection{Supportive Branch Logic}
\paragraph{Pre-registered Branch (Expected---Not Observed).} The supports\_mechanism branch is triggered when the prespecified analysis is consistent with the declared relation while planned controls do not favor a stated alternative explanation. Under this condition, the allowed conclusion is that the result supports, but does not prove, the declared relation within the design boundary. The next action requires replication under independently confirmed conditions and testing the most consequential declared boundary condition.

\subsection{Heterogeneous and Null Branch Logic}
\paragraph{Pre-registered Branch (Expected---Not Observed).} The partial\_or\_heterogeneous branch is triggered when the prespecified analysis indicates variation across declared conditions, units, or measurement contexts. The allowed conclusion is that the relation may be conditional or heterogeneous; no universal conclusion is warranted. The next action requires predefining and checking plausible moderators and improving the coverage of conditions and measurement comparability. The null\_or\_contradictory branch is triggered when the prespecified comparison does not support the declared relation or instead favors a declared alternative explanation. The allowed conclusion is that the proposed relation is not supported in this design boundary, and absence of support is not proof of absence generally. The next action requires auditing construct validity and comparison adequacy, and revising the mechanism or boundary claim before another design iteration.

\subsection{Invalid Branch and Dependency Consequence}
\paragraph{Pre-registered Branch (Expected---Not Observed).} The uninformative\_or\_invalid branch is triggered when prespecified quality-control, missingness, protocol-deviation, or validity criteria prevent interpretation. The allowed conclusion is that no scientific conclusion is warranted because the planned design did not yield interpretable evidence. The next action requires resolving the identified validity or data-quality failure before repeating the design and obtaining human confirmation of measurement, sampling, and analysis prerequisites. This branch maps to the dependency on the review ledger, which holds the canonical records for unresolved human-review requirements.

\section{Risks, Limitations, and Human Review Requirements}
\subsection{Scope Limits and Alternative Mechanisms}
\paragraph{Proof Obligation Registry PO1, PO2 (Unverified).} This proposal asserts a candidate mechanism, not a completed theorem or an observed result. The frozen formal reasoning plan remains unverified: the relaxed static equilibrium axioms, the derived mass--spin relation, and the causality-preserving step are all candidate or unverified forward-derivation steps. The counterexample analysis was not run, so no conclusion about counterexamples may be drawn, and the design must not be treated as complete until that gap is closed. The primary competing mechanism is standard stiff dense-matter equations of state combined with ordinary rotational support from binary spin-up, which already explains masses beyond the static Tolman--Oppenheimer--Volkoff limit without relaxing the nuclear equation-of-state axioms. The proposed mechanism earns its place only if it predicts a mass--spin relation that standard recycling cannot reproduce, and that relation must survive both the causality bound and the existing population of massive millisecond pulsars. Evidence from compact-binary catalogs and equation-of-state constraints sharpens the boundary: component-mass inference is sensitive to spin priors, and radius constraints depend on whether the equation of state is required to support masses above roughly two solar masses. Binary population synthesis also warns that omitted physics, such as tidal evolution, introduces systematic bias, so any numerical implementation must be declared and checked rather than assumed.~\cite{cite_3d2e6757608d5364,cite_42e341ff6471ce89,cite_81b31d0d8218f9b8}

\begin{table*}[htbp]
\centering
\footnotesize
\setlength{\tabcolsep}{3pt}
\renewcommand{\arraystretch}{1.12}
\caption{Limitation and Human-Review Decision Matrix}
\label{tab:risk-limitations-and-review-block-2}
\begin{tabular}{|p{0.23\textwidth}|p{0.23\textwidth}|p{0.23\textwidth}|p{0.23\textwidth}|}
\hline
Decision Gate & Limitation or Review Requirement & Allowed Conclusion & Release Condition \\\hline
Counterexample analysis & The counterexample batch was discarded; no counterexample was produced or excluded. & The candidate mechanism is unbounded; completeness is not established. & Qualified human review supplies or rules out a counterexample that satisfies all stated assumptions. \\\hline
Formal reasoning plan & The relaxation of the static equilibrium axioms and the derived mass--spin relation remain candidate or unverified. & No theorem status; the relation is a proposal. & Human review fixes the mathematical form of the relaxation, the mass--spin function, and the causality criterion. \\\hline
Risk gate & The canonical risk gate requires qualified human review before escalation. & The plan is not releasable as complete. & Reviewer signs off on assumptions, verification plan, and formal claim. \\\hline
Workflow degradation & LLM or workflow degradation was triggered during analysis. & Generated detail is untrusted. & Independent re-derivation or human-authored replacement of degraded outputs. \\\hline
No-information branches & Unresolved definitions (relaxed axioms, accretion support, mass--spin function, causality check) and proof obligations PO1 and PO2 are unavailable. & The dependency does not update theorem status; it is neither proof of failure nor a reason to invalidate the plan. & Resolve or review the upstream input, then re-run the dependency check. \\\hline
\end{tabular}
\end{table*}

\subsection{Concrete Human-Review Decisions and Release Conditions}
\paragraph{Proof Obligation Registry PO1, PO2 (Unverified).} Before any escalation, a qualified reviewer must make four binding decisions. First, supply the exact mathematical relaxation of the static equilibrium axioms and the functional form of the mass--spin relation, because these are referenced but undefined. Second, specify the equations governing binary mass and angular-momentum transfer, together with the formal criterion for testing a causality violation, since the speed-of-sound constraint is currently only a symbolic placeholder. Third, define the measurement methods for gravitational mass above two solar masses and for spin period, so that model outputs are comparable to the observed massive millisecond pulsar population. Fourth, adjudicate the discarded counterexample batch and the degraded workflow stage, confirming that no generated claim, formalization, or protocol detail survives without human replacement. Release of this section as complete requires all four decisions recorded and the two proof obligations either discharged or explicitly deferred; until then, the proposal stands as a scoped candidate mechanism whose conclusions are withheld, not refuted.

\subsection{Transition to the Next Stage}
The scoped premises inherited from the expected-outcomes stage are now bounded by this ledger: each prespecified branch, including the no-information branches, maps to a concrete review decision rather than to a vague caveat. The next stage should consume these decisions as inputs, advancing the formal reasoning plan only after the relaxation, the mass--spin function, and the causality criterion are fixed, and treating any outcome branch as a design consequence rather than as an observed finding.

\section{Definitions, Propositions, and Proof Obligations}
\subsection{Control Panel Scope and Symbol Domain}
This section owns the definition ledger, the candidate propositions, and the proof obligations for the mechanism-replacement route: relaxing the static equilibrium axioms of the dense-matter equation of state so that rotational and thermal support from binary accretion can carry masses beyond the static Tolman-Oppenheimer-Volkoff (TOV) limit while preserving the causal speed bound, and so that the resulting mass-spin relation separates massive millisecond pulsars from isolated magnetic-dipole spin-down pulsars. The symbol domain is fixed as follows. Gravitational mass M is measured in solar masses and spin period P in milliseconds, both taken as non-negative real quantities. The causal bound is the dimensionless constraint on the local speed of sound relative to the speed of light, and it is imposed as a hard control variable rather than a tunable parameter. The evolutionary channel is a discrete label distinguishing isolated spin-down from binary recycling. The relation R that maps P to M is a proposed object, not a closed-form result: its functional form is referenced by the frozen handoff but left undefined, and it is registered as an unresolved formal input owned by this ledger. Because the route is a proposal without observed results, every relation below is a candidate or unverified derivation target, and no numerical mass, radius, or period value is asserted as a finding.

\subsection{Definition Ledger (Owner Record)}
\paragraph{Definition.} The following entries are the complete definition ledger for this route; other sections cite it rather than restating it. Candidate-formalization symbols are usable as written. Needs-human-input symbols are placeholders whose mathematical content is not yet fixed, and they gate the proof obligations that depend on them. D1 (symbol A1, needs human input): the set of relaxed static-equilibrium axioms applied to the core mass-spin relation; its specific relaxation is undefined. D2 (symbol A2, needs human input): the binary accretion support mechanism as a continuous thermodynamic and mechanical perturbation; its governing equations are undefined. D3 (symbol M, candidate formalization): gravitational mass of the neutron star, in solar masses. D4 (symbol P, candidate formalization): spin period, in milliseconds. D5 (symbol R, needs human input): the derived mechanism mass-spin relation M = R(P); the functional form is referenced but not defined. D6 (symbol A3, candidate formalization): the evolutionary channel, isolated versus binary recycling. D7 (symbol c, needs human input): the causal speed-limit check on the speed of sound; the formal test criteria are missing. Assumption A1 (candidate formalization): the static-equilibrium axioms can be relaxed to permit mass accumulation beyond the static TOV limit via rotational and thermal support without violating the causality bound. Assumption A2 (candidate formalization): binary mass and angular-momentum transfer can be modeled as a continuous perturbation modifying the effective equilibrium conditions. Assumption A3 (candidate formalization): companion ablation acts primarily on the surface and magnetosphere, leaving the core mass-spin relation governed by the relaxed effective axioms.~\cite{cite_55957f4ce67bd972,cite_c369df652181e8f7,cite_f16cd637681fe829}

\subsection{Numbered Symbol-Domain and Threshold Convention}
M $\in$ $\mathbb{R}$\_\{>0\}, P $\in$ $\mathbb{R}$\_\{>0\}, c\_s/c $\leq$ 1, $\chi$ $\in$ \{isolated, recycled\}, M\_\{static,max\} = $\sigma$(M\_\{static,max\}) $\in$ $\mathbb{R}$\_\{>0\}, R : $\mathbb{R}$\_\{>0\} $\rightarrow$ $\mathbb{R}$\_\{>0\}, R(P) > M\_\{static,max\} $\Rightarrow$ $\chi$ = recycled

\subsection{Reading the Domain Relation}
The relation fixes the conventions that the rest of the theory cites. The mass and period are positive reals; the causality control is the dimensionless inequality c\_s/c $\leq$ 1, which must hold everywhere in the relaxed solution and is not a free parameter. The channel $\chi$ is discrete. The static maximum mass M\_\{static,max\} is treated as a positive threshold quantity whose value is supplied by the equation-of-state model rather than asserted here. The relation R is declared as a map from period to mass, and the final implication states the route's proposed classification criterion: a relaxed solution whose mass exceeds the static threshold is assigned to the recycled channel. This implication is a candidate criterion, not a proved theorem; it is exactly what the proof obligations below are required to establish or reject.

\subsection{Candidate Propositions and Lemma Registry}
\paragraph{Proposition (Candidate).} Two candidate propositions anchor the route. P1 (candidate formalization): relaxing the static-equilibrium axioms of the dense-matter equation of state with rotational and thermal support permits masses beyond the static TOV limit without violating causality. P2 (candidate formalization, the target proposition): the relaxed axioms predict a distinct mass-spin relation that separates massive neutron stars above roughly two solar masses with millisecond periods from isolated magnetic-dipole spin-down pulsars. The lemma registry decomposes these into eight units. L1 (Candidate) formalizes P1 under assumptions A1-A3. L2 (Candidate) formalizes P2 under the same assumptions. L3 (Candidate) states the proposed first derivation step: the static equilibrium constraint is modified by a continuous perturbation term tied to binary accretion support, and it depends on the undefined relaxation content of D1 and D2. L4 (Unverified) asserts existence of the relation R(P) for the relaxed axioms, gated by the undefined D5. L5 (Unverified) evaluates the mass limit, claiming R(P) yields M above the static maximum for millisecond periods. L6 (Unverified) checks that the relaxation preserves the causality bound, gated by the undefined D7. L7 (Unverified) separates the isolated magnetic-dipole track from the R(P) core-equilibrium curve. L8 (Unverified) combines the mass-limit and channel-separation steps into the final population distinction. Every lemma is a proposed, source-bounded step; none is a completed proof.

\subsection{Proof-Obligation Registry and Dependency Closure}
\paragraph{Proof Obligation Registry PO1, PO2 (Unverified).} Two proof obligations carry the route's formal burden, and both are unverified proposal work requiring human input. PO1 requires establishing that the relaxed axioms admit a causal equilibrium sequence with mass above the static TOV limit; it depends on A1-A3, the undefined relaxation content D1-D2, the causality check D7, and the derivation chain S1-S6. PO2 requires establishing that the resulting relation R(P) is distinct enough to separate the massive recycled population from the isolated spin-down population; it depends on A1, A3, the channel label D6, the undefined relation D5, and the same derivation chain. The dependency closure is the matrix below. A cell marked gate means the obligation cannot be discharged until that definition is fixed by qualified human review; a cell marked consumes means the obligation relies on that lemma or derivation step. Because D1, D2, D5, and D7 are all needs-human-input, neither obligation is currently closable, and this is a precise procedural dependency rather than evidence of failure.

\begin{table*}[htbp]
\centering
\footnotesize
\setlength{\tabcolsep}{3pt}
\renewcommand{\arraystretch}{1.12}
\caption{Dependency-Closure Matrix}
\label{tab:definitions-and-propositions-dp-dep-matrix}
\begin{tabular}{|p{0.1533\textwidth}|p{0.1533\textwidth}|p{0.1533\textwidth}|p{0.1533\textwidth}|p{0.1533\textwidth}|p{0.1533\textwidth}|}
\hline
Unit & Role & PO1 & PO2 & Status & Gate \\\hline
A1-A3 & premises & consumes & consumes & candidate & no \\\hline
D1, D2 & relaxation content & gate & gate & needs human input & yes \\\hline
D5 & relation R(P) & consumes & gate & needs human input & yes \\\hline
D7 & causality check & gate & gate & needs human input & yes \\\hline
S1-S6 & derivation chain & consumes & consumes & unverified & no \\\hline
L1-L8 & lemma registry & consumes & consumes & candidate/unverified & no \\\hline
\end{tabular}
\end{table*}

\subsection{No-Information Decision Branches}
\begin{itemize}
\item The control panel registers six no-information branches, one per unresolved gate and one per proof obligation. When a gate is unavailable, the corresponding branch withholds any theorem-status update; it is neither a proof of failure nor a reason to invalidate the research plan. The branches are: No-information: D1, No-information: D2, No-information: D5, No-information: D7, No-information: PO1, and No-information: PO2. Each maps to the same next action, resolve or review the upstream input. The concrete consequence is that the causality claim (L6) and the existence of the relation (L4) remain open until D7 and D5 are fixed, and the relaxation content (L3) remains open until D1 and D2 are fixed. These are procedural dependencies owned by this ledger; the detailed human-confirmation and release criteria are handled by the review section, which this section references without restating.
\end{itemize}

\subsection{Evidence-Backed Context and Boundary}
The background that motivates the relaxed-axiom route is established upstream. The observed neutron-star radius and tidal constraints from the binary merger GW170817 are sensitive to whether the analysis additionally requires the equation of state to support masses above roughly 1.97 solar masses, which is the empirical anchor for the massive-neutron-star regime this route addresses. Joint microscopic and macroscopic constraints from heavy-ion collision data improve the dense-matter pressure inference, sharpening the causal stiffness bound that the relaxation must respect. The binary-accretion perturbation in A2 is grounded in the documented role of tides, mass transfer, and mergers in setting rotation rates, and in the angular-momentum balance between accretion spin-up and wind braking. Two boundary facts constrain the design. First, the standard population-synthesis machinery that supplies accretion and tidal prescriptions explicitly excludes common-envelope and contact-binary phases, so the continuous-perturbation idealization of A2 is bounded away from those regimes. Second, the same machinery treats mass-transfer efficiency and magnetic braking as the dominant sources of uncertainty, which is why D2 is registered as needing formal specification rather than being asserted. These are comparison and design conclusions, not new results.~\cite{cite_1e5d552669fc0c56,cite_2908f4a456e2afd6,cite_42e341ff6471ce89,cite_6c11517977d0b39d,cite_c369df652181e8f7}

\subsection{Failure Conditions and Transition}
The route's declared failure boundary is explicit and is what the derivation and counterexample stage must interrogate. The claim fails if the relaxed static-equilibrium axioms violate the causality bound, if the predicted mass-spin relation contradicts the observed population of massive millisecond pulsars, or if standard rotational support alone fully explains the observations without relaxing the static axioms. The third condition is the live alternative: a stiff equation of state combined with ordinary binary spin-up may already account for the high masses, in which case the relaxation adds no explanatory work. A counterexample that defeats the route must satisfy all assumptions A1-A3 and the causality control simultaneously; a violation that only holds outside the declared domain does not qualify. The counterexample analysis for this route was not completed and is deferred to qualified human review, so no counterexample is claimed here. This section closes the formal problem and hands the next stage a fixed symbol domain, the candidate propositions P1 and P2, the lemma registry L1-L8, the proof obligations PO1 and PO2, and the precise gates D1, D2, D5, and D7 that must be resolved before any theorem status can change.

\section{Forward Derivation and Counterexample Search Plan}
\subsection{Consumed Premises and Derivation Scope}
This section advances the argument by converting the candidate premises supplied by the definitions and propositions stage into a forward derivation obligation, a failure condition, and a counterexample decision protocol. The three formal assumptions inherited from the definition ledger---A1, the relaxation of static equilibrium axioms under rotational and thermal support; A2, the treatment of binary mass and angular momentum transfer as a continuous mechanical perturbation; and A3, the confinement of companion ablation to the surface and magnetosphere---constitute the admissible premise set. The derivation does not assert that these premises are established. It specifies what must follow if the candidate mechanism is internally coherent, and what would falsify the candidate if it is not. The target is the mass-spin relation M = R(P) and the claim that this relation distinguishes massive millisecond pulsars from isolated magnetic-dipole spin-down pulsars. No observed results are claimed; every formal step is a proposed obligation subject to human review.

\subsection{Derivation Notation and Admissible Comparators}
\paragraph{Definition.} The derivation uses the symbols fixed by the definition ledger. M denotes gravitational mass, P denotes spin period, R(P) denotes the candidate mass-spin function, and c\_s <= c denotes the causal speed-of-sound constraint. The static maximum mass M\_static\_max is the Tolman-Oppenheimer-Volkoff limit computed under unrelaxed equilibrium axioms. The admissible comparator for the mechanism claim is the standard rotational-support model: a stiff equation of state combined with centrifugal support from binary spin-up, without any relaxation of the fundamental nuclear equation-of-state axioms. This comparator is not a strawman; it is the declared alternative explanation that the derivation must distinguish from the relaxed-axiom mechanism. A second admissible comparator is the isolated magnetic-dipole spin-down track, which serves as the population boundary against which the R(P) curve is tested. Any counterexample that violates the causal bound or falls outside the millisecond-period regime is an out-of-domain boundary case, not a valid refutation of the scoped claim.

\subsection{Numbered Derivation Relation}
M = R(P; A1, A2, A3) subject to c\_s <= c and M > M\_static\_max for P in the millisecond regime

\subsection{Equation Interpretation and Lemma Chain}
The relation above is the target of the forward derivation, not a completed proof. It states that the relaxed-axiom mechanism predicts a mass-spin curve R(P) that exceeds the static maximum mass in the millisecond period range while preserving the causal bound. The derivation proceeds through a lemma chain. Lemma L3 applies the relaxation of static equilibrium axioms under the continuous perturbation of binary accretion, yielding a modified effective equilibrium condition. Lemma L4 derives the existence of the mass-spin function R(P) from that modified condition. Lemma L5 evaluates the mass limit and asserts that R(P) yields M > M\_static\_max for millisecond periods. Lemma L6 checks that the relaxation preserves the causal speed-of-sound constraint. Lemma L7 characterizes the isolated magnetic-dipole spin-down track as distinct from the R(P) curve. Lemma L8 combines the mass-limit evaluation and the evolutionary-channel distinction to separate the massive millisecond population from the isolated pulsar population. Each lemma is a proposed obligation, not a verified result. The chain is conditional on the three assumptions and on the formal definitions of the relaxation and the accretion perturbation, both of which currently require human input.

\subsection{Candidate Lemma L4 and Proof Obligations PO1, PO2}
\paragraph{Lemma Registry L4, PO1, PO2 (Unverified).} Lemma L4 (Unverified): Given the relaxed equilibrium condition from Lemma L3 and the definitions of gravitational mass M and spin period P, there exists a function R(P) such that the equilibrium sequence of the perturbed star is parameterized by P. This lemma is the load-bearing step: the entire mass-spin prediction rests on the existence and form of R(P). Its status is unverified because the specific functional form of the mechanism\_mass\_spin\_relation is referenced in the upstream plan but not defined, and the equations governing binary mass and angular momentum transfer lack formal specification. Two proof obligations gate this lemma. PO1 requires a formal demonstration that the relaxation axioms and the derived relation consistently produce massive neutron star masses within the declared validity domain without violating causality. PO2 requires a demonstration that the derived relation distinguishes the massive millisecond population from the isolated spin-down population under the admissible comparator. Both obligations are currently in a needs\_human\_input state. A failure of PO1 invalidates the mechanism claim within its scope. A failure of PO2 does not invalidate the mechanism but eliminates its discriminating power, reducing it to a description compatible with the standard rotational-support alternative.

\begin{table*}[htbp]
\centering
\footnotesize
\setlength{\tabcolsep}{3pt}
\renewcommand{\arraystretch}{1.12}
\caption{Counterexample Decision Matrix}
\label{tab:forward-derivation-and-counterexamples-fd-06}
\begin{tabular}{|p{0.1533\textwidth}|p{0.1533\textwidth}|p{0.1533\textwidth}|p{0.1533\textwidth}|p{0.1533\textwidth}|p{0.1533\textwidth}|}
\hline
Case & Satisfies A1, A2, A3? & Satisfies c\_s <= c? & Distinguishes from standard rotational support? & Classification & Response \\\hline
Causality-violating R(P) at millisecond P & Yes & No & N/A & Would-falsify & Reject mechanism claim; revise A1 or A2 \\\hline
R(P) matches standard rotational-support curve & Yes & Yes & No & Would-falsify (discrimination) & Retain mechanism as compatible but non-discriminating; revise A3 or scope \\\hline
R(P) contradicts observed massive MSP population & Yes & Yes & Yes & Would-falsify (empirical) & Reject mechanism claim within declared scope \\\hline
Relaxation applied to non-millisecond or isolated systems & Yes & Yes & N/A & Out-of-domain boundary & No conclusion; outside declared validity domain \\\hline
Counterexample assumes ablation alters core equilibrium & No (violates A3) & N/A & N/A & Out-of-domain boundary & Invalid counterexample; A3 is an admissible premise \\\hline
Missing formal definition of R(P) or accretion equations & Indeterminate & Indeterminate & Indeterminate & No-information & Withhold theorem status update; resolve upstream input via human review \\\hline
\end{tabular}
\end{table*}

\subsection{Failure Conditions, Limitations, and Transition}
The declared failure condition is precise: the mechanism claim fails if the relaxed static equilibrium axioms violate the causality bound, if the predicted mass-spin relation contradicts the observed population of massive millisecond pulsars, or if standard rotational support models fully explain the observations without requiring the relaxed axioms. The third condition is the most consequential for the argument's contribution, because the standard rotational-support comparator is already a viable explanation for massive millisecond pulsars in the literature. The derivation must therefore produce a discriminating prediction, not merely a consistent one. A limitation constrains this section's force: the counterexample analysis stage was discarded due to an unavailable response, so no generated counterexample, formalization, or protocol detail from that batch is retained. The decision matrix above is a proposed search plan, not a completed audit. The verification plan and the counterexample-or-boundary analysis both require human input before this design can be treated as complete. The transition to the expected outcomes stage is concrete: the outcome branches map the prespecified analysis to the lemma and proof-obligation states defined here. A no-information branch, such as the unresolved definition of R(P), does not update theorem status and does not invalidate the research plan; it specifies the dependency that must be resolved before the outcome branches carry evidential weight.

\section{Computational Evidence (Numerical Simulation; Non-empirical)}
Q1 (final v0) answers: Can the joint distribution of braking index n and characteristic age tau\_c statistically distinguish isolated magnetic-dipole evolution from binary recycling populations, given specific torque model constraints? Model family: Explicit Runge-Kutta 4th Order. Execution mode: NUMERICAL\_SIMULATION. Result kind: SIMULATED. Empirical claim status: NOT\_EMPIRICAL. Model-internal hypothesis relation: CONSTRAINED. Model-internal result summary: Within the audited reduced ODE, the isolated track increases from P=0.500 s to 0.662 s, whereas the recycling track reaches P\_min=0.226 s during accretion and ends at 0.281 s after resumed dipole spin-down. This establishes model-internal directional separation, but does not by itself validate population-level n-tau\_c discrimination because the executable MathIR contains two deterministic tracks rather than the documented Monte Carlo outer loop. Numerical quality: NOT\_REPORTED. Applicability conditions: Stellar geometry, electromagnetic constants, and moment of inertia are absorbed into effective braking constants K\_dipole and K\_acc.; Magnetic inclination evolves smoothly via exponential decay without discontinuous glitches or magnetospheric state transitions.; Accretion torque is active exclusively during t < t\_acc with nonzero Mdot; isolated dipole torque applies otherwise.. Limitations: Excludes glitch events, magnetospheric state transitions, and variable moment of inertia.; Assumes fixed stellar geometry and absorbs unresolved electromagnetic factors into effective constants.; Observational selection effects and survey detection thresholds are not dynamically coupled to the simulation loop.. Iteration lineage: v0: CONSTRAINED (Within the audited reduced ODE, the isolated track increases from P=0.500 s to 0.662 s, whereas the recycling track reaches P\_min=0.226 s during accretion and ends at 0.281 s after resumed dipole spin-down. This establishes model-internal directional separation, but does not by itself validate population-level n-tau\_c discrimination because the executable MathIR contains two deterministic tracks rather than the documented Monte Carlo outer loop.). Supplementary mathematical-model PDF: quantitative\_mathematical\_models.pdf\#Q1.

Q2 (final v1) answers: How does the radial thermal evolution of a massive neutron star differ when spin-down-associated core compression triggers a localized phase-transition heat source, compared with standard cooling and with a superfluid-modified cooling scenario? Model family: spherical\_radial\_thermal. Execution mode: NUMERICAL\_SIMULATION. Result kind: SIMULATED. Empirical claim status: NOT\_EMPIRICAL. Model-internal hypothesis relation: INCONCLUSIVE. Model-internal result summary: Q2 v1 baseline spherical-radial thermal PDE completed with finite output, explicit stability, spherical-origin regularity, and boundary checks passing; this baseline-only run does not establish scenario separation or empirical pulsar validation. Numerical quality: NUMERICALLY\_VERIFIED. Applicability conditions: The star is spherically symmetric and the executable PDE uses the canonical coordinate x as the physical radial coordinate r normalized by the stellar radius.; Spin-down induced core compression and phase-transition onset are represented by a prescribed localized latent-heat source rather than dynamically solved spin, density, or moment-of-inertia equations.; Thermal microphysics is reduced to scaled constant transport, heat-capacity, and source coefficients.; The surface boundary is approximated by a fixed scaled surface temperature representing atmosphere and radiative-transfer effects.; Units are normalized so that the temperature field is scaled by a reference core temperature, x=1 corresponds to the stellar radius, and t=1 corresponds to a characteristic cooling time.. Limitations: The model does not solve spin evolution, moment of inertia, or central density dynamically; spin-down compression is parameterized by a prescribed source.; Microphysical equation-of-state effects, latent heat, neutrino emissivity, and superfluid suppression are represented by scaled coefficients rather than density-temperature tables.; The surface boundary is simplified; atmosphere radiative transfer and gravitational redshift are not solved.; Spherical symmetry and a uniform radial grid exclude multidimensional effects and sharp phase-boundary dynamics.; No observational data are ingested; comparison to X-ray cooling curves requires external calibration and uncertainty treatment.. Iteration lineage: v0: SUPPORTED\_WITHIN\_MODEL (The audited reduced ODE crosses rho\_crit=5.12e17 kg/m\textasciicircum{}3 at approximately 9.85e11 s in all spin-down cases. At 3*tau\_cool, the standard, phase-transition, and phase-transition-plus-superfluid scenarios yield surface temperatures of about 4.98e4 K, 5.32e4 K, and 2.36e5 K, respectively; phase\_fraction reaches 1 only in the two transition scenarios. The mechanism therefore produces distinguishable model-internal thermal tracks, conditional on the effective cooling, latent-heat, and compression proxies.); v1: INCONCLUSIVE (Q2 v1 baseline spherical-radial thermal PDE completed with finite output, explicit stability, spherical-origin regularity, and boundary checks passing; this baseline-only run does not establish scenario separation or empirical pulsar validation.). Supplementary mathematical-model PDF: quantitative\_mathematical\_models.pdf\#Q2.

\appendices
\section{Energy-Condition Taxonomy, Symbols, and Boundary Defense}
\subsection{Scope and Purpose}
This appendix establishes the boundary conditions for the proposed mechanism, which relaxes static equilibrium axioms to permit mass accumulation beyond the static Tolman-Oppenheimer-Volkoff limit. The argument relies on a strict separation between causal speed limits and gravitational energy conditions. While the proposal assumes the speed of sound remains subluminal (c\_s <= c), this constraint is distinct from the Null Energy Condition (NEC) or the Strong Energy Condition (SEC). The following taxonomy prevents the conflation of these independent physical requirements.

\begin{table*}[htbp]
\centering
\footnotesize
\setlength{\tabcolsep}{3pt}
\renewcommand{\arraystretch}{1.12}
\caption{Energy Condition and Causality Taxonomy}
\label{tab:appendix-variables-and-definitions-taxonomy-table}
\begin{tabular}{|p{0.23\textwidth}|p{0.23\textwidth}|p{0.23\textwidth}|p{0.23\textwidth}|}
\hline
Condition & Mathematical Form & Role in Proposal & Logical Boundary \\\hline
**Causality Bound** & c\_s <= c & Independent Constraint & Preserved by assumption; not derived from energy conditions. \\\hline
**NEC** & T\_\{ab\}k\textasciicircum{}a k\textasciicircum{}b >= 0 & Background Assumption & Required for standard stability; does not imply SEC. \\\hline
**SEC** & (T\_\{ab\} - 1/2 T g\_\{ab\})k\textasciicircum{}a k\textasciicircum{}b >= 0 & Not Assumed & Violation permitted in rotating frames; not a substitute for NEC. \\\hline
**AANEC** & integral of T\_\{ab\}k\textasciicircum{}a k\textasciicircum{}b >= 0 & Not Required & Does not establish trapped surfaces in this context. \\\hline
\end{tabular}
\end{table*}

\subsection{Causality and Equilibrium Separation}
The relaxation of equilibrium axioms introduces rotational and thermal support terms that modify the effective mass limit. This modification must not violate the causal speed limit, defined as the condition that the speed of sound in dense matter remains below the speed of light. The taxonomy above clarifies that satisfying the causality bound (c\_s <= c) does not automatically satisfy the Strong Energy Condition (SEC), nor does it guarantee the validity of the Null Energy Condition (NEC) under all perturbations. The proposal treats the causality bound as an independent constraint (V5) that must be explicitly verified against the relaxed axioms (V1), rather than assuming it follows from standard static energy conditions.

\subsection{Unresolved Operationalizations}
\begin{itemize}
\item The following variables require formal definition to complete the boundary defense:
\item **V1 (Relaxed Axioms):** The specific mathematical relaxation of static equilibrium axioms is currently undefined.
\item **V2 (Mass Measurement):** The operational method for measuring masses >2.0 M\_sun (e.g., Shapiro delay) is not specified.
\item **V3 (Spin Measurement):** The measurement protocol for spin period is missing.
\item **V4 (Binary Transfer):** The equations governing binary mass and angular momentum transfer require formal specification.
\item **V5 (Causality Test):** The formal criteria for testing causality violation in the relaxed model are absent.
\end{itemize}

\section{Idea Source Checkpoints and Direction Selection Audit}
\subsection{Checkpoint Status and Directional Anchor}
The proposal rests on a single selected mechanism, direction mechanism\_replacement, whose central hypothesis asserts that relaxing the static equilibrium axioms of the dense matter equation of state to incorporate rotational support from binary mass accretion formally unifies isolated spin-down and binary evolution. Three available audit snapshots register this direction without establishing a temporal order between them. The first snapshot records the direction as a candidate with claim scope candidate\_mechanism and target gap 'Required slot candidate\_mechanism remains uncovered in relevant evidence cluster Massive Neutron Star Masses and Dense Matter Equations of State.' The second snapshot preserves the same hypothesis under a divergent title, Dynamical Inclination and Wind Braking in Isolated Pulsar Spin-Down: Deriving the Effective Braking Index and Validity Domain, signaling that the portfolio audit did not resolve the naming of the selected direction. The third snapshot expands the hypothesis into a full structured entry carrying the three formal assumptions, the boundary failure condition, and the discriminating observation that massive neutron stars above two solar masses must exhibit millisecond spin periods indicative of rotational support. These checkpoints are audit records, not empirical results, and they do not confirm the mechanism.

\subsection{Scope Corrections and Retained Decisions}
The retained decision preserves the mechanism as a candidate formalization rather than a verified result. The survey gap record is the load-bearing premise: the relevant evidence cluster on massive neutron star masses and dense matter equations of state leaves the candidate\_mechanism slot uncovered, so the proposal fills a structural void rather than resolving a demonstrated inconsistency. The frozen evidence supports the boundary conditions of that void. The unified dense matter equation of state treats the crust under a ground-state composition and the liquid core under a minimal npe$\mu$ composition, which is the static baseline the proposal claims to relax. GW170817 radius constraints depend on whether the analysis additionally requires the equation of state to support neutron stars above 1.97 solar masses, the electromagnetic benchmark for massive remnants. Joint microscopic and macroscopic analyses improve dense matter constraints by raising pressure, and component mass inference is strongly sensitive to spin priors. These findings bound the proposal's domain; they do not validate the relaxation. The counterexample analyzer batch was discarded after an invalid response, so no counterexample, formalization, or protocol detail from that batch is retained, and the design is not treated as complete until qualified human review replaces it.~\cite{cite_55957f4ce67bd972,cite_42e341ff6471ce89}

\subsection{Open Questions and Transition to Formal Work}
Three open questions survive the checkpoint audit and gate the next stage. The exact mathematical relaxation conditions for the static equilibrium axioms remain undefined, which is the premise that assumption A1 states but does not prove. The specific numerical methods or analytical techniques for verifying the mass-spin relation are unspecified, leaving the derived relation R(P) as a placeholder rather than a computable object. The causal speed-limit check, the formal criteria for testing whether the relaxation violates the c\_s <= c bound, is likewise missing. These gaps are procedural dependencies on the definition ledger and the proof-obligation registry, not defects of the idea itself. The transition to the derivation and falsification stage is therefore concrete: the next section must consume the candidate lemmas and proof obligations to turn the relaxation premise and the mass-spin relation into a numbered equation chain, then map the prespecified outcome branches, including the no-information branches, to the relevant lemma or proof obligation. Until the relaxation conditions and the causality test are supplied, the proposal remains a candidate mechanism bounded by the survey's uncovered slot.

\section{Evidence Coverage, Unknown Items, and Review Checklist}
\subsection{Evidence Coverage and Logical Boundaries}
The proposal's mechanism is grounded in the intersection of binary evolution dynamics and dense matter constraints. Binary interactions, including tides, mass transfer, and mergers, are established mechanisms that alter stellar rotation rates and angular momentum distributions (W2001595193). However, modeling these interactions requires careful boundary conditions; omitting tidal evolution introduces systematic biases in population synthesis (W2059879627), and current stellar evolution modules explicitly exclude common envelope phases (W2095725842). On the microphysical side, equation of state constraints are tightly coupled to observational data. Heavy-ion collision data improves constraints on dense matter pressure (W3182570346), while gravitational wave observations provide radius limits that depend on assumptions about maximum neutron star masses (W2807309204). The logical boundary for this proposal is the separation of standard rotational support from the proposed relaxation of static equilibrium axioms. While binary spin-up is a known phenomenon, the specific claim that static equilibrium axioms themselves can be relaxed without violating causality remains a candidate mechanism requiring formal proof.~\cite{cite_1e5d552669fc0c56,cite_3d2e6757608d5364,cite_c369df652181e8f7}

\subsection{Unknown Items and Formal Dependencies}
Two critical unknowns prevent the immediate verification of the proposed mechanism. First, the specific functional form of the 'mechanism\_mass\_spin\_relation' is referenced but not defined (unknown-14). Without this mathematical specification, the derivation steps linking relaxed axioms to observable mass-spin tracks remain incomplete. Second, the counterexample analysis stage was degraded due to an unavailable LLM response (unknown-4). Consequently, no generated claims or formal counterexamples from that batch are retained. This absence means that potential violations of the causality bound or tidal deformability constraints have not been formally tested against the relaxed axioms. These items are not merely missing data but represent structural gaps in the formal reasoning plan. The proposal cannot advance to empirical testing until the mass-spin relation is explicitly formulated and the counterexample analysis is completed via qualified human review.

\subsection{Release Criteria for Future Claims}
\paragraph{Human-review checklist.} The release of any claims regarding the unified formation mechanism depends on resolving the formal dependencies identified in the theory spine. The proof obligations PO1 and PO2 are currently marked as needs\_human\_input. Release criteria require: (1) A formal definition of the mass-spin relation R(P) that satisfies the causality bound c\_s <= c under the relaxed axioms. (2) A completed counterexample analysis confirming that the relaxed axioms do not inevitably violate tidal deformability constraints or standard rotational support models. (3) Verification that the predicted relation distinguishes massive millisecond pulsars from isolated spin-down pulsars without requiring ad hoc adjustments to the equation of state. Until these criteria are met, the mechanism remains a candidate hypothesis with unverified formal reasoning.\bibliographystyle{IEEEtran}
\bibliography{references}
\end{document}
